\section{问题二的模型建立与求解}

\subsection{求解思路}

先在 B1 上估计经典标度律，并在任何结论之前检验 B1 的损失列是否含有真实训练噪声；再用独立
的观测表做外部验证；随后在半合成表上拟合包含质量的候选形式，按预先设定的采纳条件决定
是否把它作为广义标度律；最后在候选形式下推导边际效用、弹性与质量–规模替代条件。

\subsection{经典标度律的估计}

以 $N$（参数个数）与 $D$（Token 数）的原始计数拟合式 \eqref{eq:classic}。沿用
文献\cite{hoffmann2022chinchilla}的做法，对 $\log L$ 的残差取 Huber 损失
\cite{huber1964robust}（$\delta=10^{-3}$），并以 L-BFGS-B 算法\cite{byrd1995lbfgsb}
从 108 个初值出发求全局最优。不采用对数线性最小二乘，因为它隐含 $E=0$，会把不可约损失
吸收进指数。

B1 由同一模型族的 8 条训练轨迹组成，对应 Pythia 系列的 8 个规模\cite{biderman2023pythia}，
共 1176 个检查点。同一轨迹上的检查点高度相关，因此不确定性以轨迹为重抽样单位
（假设 4）。表\ref{tab:q2-params}并列报告三类诊断：以轨迹为单位的自助法区间、逐一删去
一条轨迹后的参数范围，以及剖面目标函数——把某一参数固定为估计值的
$1\pm2\%$ 并重新优化其余参数后，目标函数相对最小值的倍数。

\begin{table}[htbp]
  \centering
  \caption{经典标度律参数及三类不确定性诊断}
  \label{tab:q2-params}
  \begin{tabular}{lrrrr}
    \toprule
    参数 & 估计值 & 自助法 95\% 区间 & 删一轨迹范围 & 剖面倍数（$-2\%$/$+2\%$） \\
    \midrule
    $E$ & 1.68982 & [1.6897, 1.6899] & [1.68982, 1.68984] & 346.7 / 385.4 \\
    $A$ & 406.268 & [405.477, 406.776] & [406.154, 406.636] & 2.953 / 2.855 \\
    $\alpha$ & 0.339983 & [0.339867, 0.340051] & [0.339966, 0.340033] & 65.66 / 65.03 \\
    $B$ & 409.744 & [409.393, 410.090] & [409.615, 409.891] & 5.024 / 4.795 \\
    $\beta$ & 0.279888 & [0.279846, 0.279930] & [0.279871, 0.279906] & 80.14 / 78.67 \\
    \bottomrule
  \end{tabular}
\end{table}

三类诊断回答不同的问题，不能互相替代。删一轨迹的范围在全部 5 个参数上都窄于自助法区间：
自助法有放回地抽取 8 条轨迹，单次重抽样可能重复或遗漏多条轨迹，扰动远强于只删去一条；
任何单条轨迹对参数的影响都不超过其估计值的 0.0905\%。剖面倍数在偏离 2\% 时均明显大于 1，
说明 5 个参数都被数据确定。

\subsection{主拟合数据的生成公式识别}

上述区间极窄，这正是需要先解释的现象。把文献\cite{hoffmann2022chinchilla}发表的常数
（$E=1.69$，$A=406.4$，$\alpha=0.34$，$B=410.7$，$\beta=0.28$）不经拟合直接代入 B1，
其中位绝对相对误差为 $3.24\times10^{-5}$，而 B1 以 4 位小数存储所对应的舍入量级为
$2.11\times10^{-5}$，二者之比仅 1.54，没有为真实训练噪声留下余地。作为对照，同样的检验在
半合成表 B2 上相差 21413 倍，说明该检验能区分不同性质的数据，而非对任何表都给出同一结论。

因此 B1 的损失列与“按已发表公式逐点计算”一致。在 B1 上拟合经典标度律，证明的是估计方法
能够恢复生成公式——重新拟合得到的参数与发表值的差异不超过 0.24\%——而\textbf{不是}对真实
模型规模规律的新发现，表\ref{tab:q2-params}中的窄区间也不代表对真实规律的精确认识。
关于真实模型的经验证据只能来自下面的独立验证。

\subsection{外部验证与有效域}

\begin{table}[htbp]
  \centering
  \caption{经典标度律的外部验证}
  \label{tab:q2-validation}
  \begin{tabular}{llrrl}
    \toprule
    数据 & 性质 & 行数 & 中位绝对相对误差 & 说明 \\
    \midrule
    B5 文献数据 & 观测 & 44 & 7.52\% & 独立的模型族与机构 \\
    B4 剔除同族行 & 观测 & 49 & 8.07\% & 剔除后才构成盲验证 \\
    B4 含同族行 & 观测 & 57 & 8.62\% & 含信息泄漏，仅作对照 \\
    B2 & 半合成 & 1029 & 32.15\% & 校准数据，非观测证据 \\
    \bottomrule
  \end{tabular}
\end{table}

在独立的观测数据上，经典标度律的中位相对误差约为 8\%（表\ref{tab:q2-validation}）。
模型的有效域为 B1 覆盖的范围：$N\in[7.054\times10^{7},\,1.197\times10^{10}]$，
$D\in[1.340\times10^{8},\,2.999\times10^{11}]$；在此之外的计算均属外推。赛题要求结合
B9、B10 讨论百亿参数以上的外推：B9 只给出大模型的规格而无损失列，B10 是组织方提供的模型
输出，以 B10 拟合再以 B10 检验会构成循环论证，二者都不能检验外推的正确性。因此本文对
有效域以外、包括百亿参数以上的预测一律按外推表述，并注明其距有效域的距离。

\subsection{包含数据质量的候选形式及其采纳检验}

以“质量改变数据的有效规模”为出发点，考虑候选形式
\begin{equation}
  L(N,D,Q)=E+A N^{-\alpha}+B\left(Q^{\gamma}D\right)^{-\beta},\qquad \gamma>0,
  \label{eq:candidate}
\end{equation}
当 $Q=1$ 时它退化为经典形式。可同时观测 $N,D,Q$ 与损失的只有半合成表：B6 的 360 行
完全包含于 B7 的 450 行（共享格点上损失完全一致），故只用 B7；B8 的质量方向与 B7 在全部
150 个审计格点上相反，其语义在现有材料中无法确定，予以隔离，不参与任何拟合。

在 B7 上拟合式 \eqref{eq:candidate} 得 $\hat\gamma=1.189$，$\hat\beta=0.0699$，
二者之积 $\hat\gamma\hat\beta=0.0831$。逐一留出一个质量水平的检验中，候选形式的
$\log L$ 预测误差比不含质量的基线降低 37.6\%。这些数值只说明候选形式在半合成表内部的
近似能力，属于半合成校准结果，不是关于真实模型的观测结论。

采纳条件在拟合前固定，包括参数取值、方向稳定、内部有效、可识别性、与真实 1M 数据的一致性
和与经典规律的兼容性，且要求候选的质量轴与问题一的质量评分之间存在第一方映射。结果为：
\begin{enumerate}
  \item \textbf{可识别性不通过。}在 $\gamma$ 的剖面上，把 $\gamma$ 改为估计值的 0.9 倍或
        1.1 倍，目标函数只增大 0.8\% 与 0.5\%；$\gamma$ 变化 4 倍时乘积 $\gamma\beta$ 只
        变化 1.29 倍，数据约束的是乘积而不是 $\gamma$ 本身；
  \item \textbf{与真实 1M 数据的一致性不通过。}问题一的配比质量 $Q(\boldsymbol{p})$ 与
        1M 宏平均损失的 Spearman 相关为 $-0.025$，按已映射份额分为三档后分别为 0.056、
        0.024、$-0.093$，没有稳定的“质量越高损失越低”关系；
  \item \textbf{与经典规律的兼容性不通过。}候选形式的 $N$--$D$ 部分在剔除同族行的 B4 上
        中位误差为 9.88\%，超过经典规律的 8.07\%；且质量项只在 $D\ge10^{10}$ 的范围内被识别；
  \item \textbf{不存在第一方映射。}B7 的质量取值是半合成设计变量，而问题一的 $Q$ 是以 A1
        为参照的相对秩评分，二者单调相关或取值范围相近都不能证明尺度等价。
\end{enumerate}
因此本文\textbf{不采纳}包含质量的广义标度律。可用于后续问题的标度律是经典形式
\eqref{eq:classic}，其中不含质量维度；问题一的质量评分与 1M 配比响应面作为独立的
信息保留，不与标度律合成为跨规模的质量规律。配比 $\boldsymbol{p}$ 同样不进入标度律：
$Q(\boldsymbol{p})$ 不是配比的充分概括，而配比响应面只在 1M 上有效。

\subsection{边际效用、弹性与质量–规模替代条件}

以下结论是式 \eqref{eq:candidate} 的解析推论，对任意正参数成立，不依赖任何数值估计；
它们回答“若质量以该形式进入标度律，会有什么结构”，而不构成经验结论。

\textbf{边际效用。}
\begin{equation}
  \frac{\partial L}{\partial N}=-\alpha A N^{-\alpha-1}<0,\quad
  \frac{\partial L}{\partial D}=-\beta B Q^{-\gamma\beta}D^{-\beta-1}<0,\quad
  \frac{\partial L}{\partial Q}=-\gamma\beta B Q^{-\gamma\beta-1}D^{-\beta}<0 .
\end{equation}

\textbf{弹性。}记 $\varepsilon_X=\partial\ln L/\partial\ln X$，则
\begin{equation}
  \varepsilon_N=-\frac{\alpha A N^{-\alpha}}{L},\qquad
  \varepsilon_D=-\frac{\beta B (Q^{\gamma}D)^{-\beta}}{L},\qquad
  \varepsilon_Q=\gamma\,\varepsilon_D .
\end{equation}
质量弹性恰为数据弹性的 $\gamma$ 倍：在该形式下，质量的作用完全通过数据项实现，数据项在
损失中占比越小，提升质量的收益越小。

\textbf{替代条件。}沿等损失曲线 $\mathrm{d}N/\mathrm{d}Q=-L_Q/L_N<0$，即提高质量可减少
达到同一损失所需的参数量。把质量由 $Q_0$ 提升到 $Q_1$ 使数据项下降
\begin{equation}
  R=B D^{-\beta}\left(Q_0^{-\gamma\beta}-Q_1^{-\gamma\beta}\right)>0 .
\end{equation}
若要在质量不变的情况下只靠增加参数获得同样的损失下降，所需参数量 $N'$ 满足
$A N'^{-\alpha}=A N^{-\alpha}-R$，即
\begin{equation}
  N'=\left(N^{-\alpha}-R/A\right)^{-1/\alpha},\qquad
  \Delta N=N'-N .
  \label{eq:equiv}
\end{equation}
式 \eqref{eq:equiv} 有限当且仅当 $R<A N^{-\alpha}$：质量提升带来的收益一旦超过参数项
剩余的全部可下降空间，就不存在与之等价的参数增量。“质量提升 0.1 等价于参数增加多少”
由此可在给定 $(N,D,Q_0)$ 下直接计算；但由于式 \eqref{eq:candidate} 未被采纳，本文不把
任何具体数值作为真实模型的结论。单位算力下两条路径的比较还依赖问题三的成本函数。
